\section{Vstavané systémy — podrobný prehľad}
\label{sec:embedded-detailed}

Táto kapitola rozpracováva kľúčové aspekty návrhu vstavaných systémov. Je určená pre inžinierov a študentov, ktorí potrebujú praktické rady pri výbere komponentov, návrhu napájania, riadení časovania a správe softvéru.

\subsection{Architektúra systému}
Základom návrhu je pochopenie funkčných blokov: senzory, spracovanie, pamäť, komunikácia a akčné členy. Pri menších systémoch postačuje jednoduchá jednovrstvová architektúra, pri zložitejších sa uplatní viacvrstvový prístup s jasným oddelením zberu dát, lokálneho spracovania a odosielania do cloudu.

\begin{figure}[h]
\centering
\begin{tikzpicture}[
node distance=11mm and 12mm,
block/.style={draw, rounded corners, minimum width=3.1cm, minimum height=0.9cm, align=center},
arrow/.style={-{Stealth[length=2.2mm]}, thick}
]
\node[block] (sensor) {Senzory};
\node[block, right=of sensor] (mcu) {MCU};
\node[block, right=of mcu] (act) {Akčné členy};
\node[block, below=of mcu] (comm) {Komunikácia};
\node[block, above=of comm] (cloud) {Cloud / UI / API};
\draw[arrow] (sensor) -- node[above, font=\small]{vstupné dáta} (mcu);
\draw[arrow] (mcu) -- node[above, font=\small]{telemetria} (comm);
\draw[arrow] (comm) -- node[right, font=\small]{prenos} (cloud);
\draw[arrow] (mcu) -- node[right, font=\small]{riadenie} (act);
\draw[arrow] (cloud) |- node[pos=0.25, right, font=\small]{príkazy} (mcu);
\end{tikzpicture}
\caption{Všeobecná bloková schéma vstavaného IoT zariadenia.}
\end{figure}

Pri návrhu treba zvážiť:
\begin{itemize}
\item požadovaný výpočtový výkon a spotrebu,
\item dostupnú pamäť (RAM a flash),
\end{itemize}

\subsection{Architektonický návrh a rozdelenie zodpovedností}
Dobrá architektúra oddelí zber dát, lokálne spracovanie, komunikačnú vrstvu a aplikačnú logiku. Pre návrh je užitočné formalizovať požiadavky: latenciu, spoľahlivosť (SLA), energetické limity, bezpečnostné ciele a podmienky prevádzky. Na základe týchto požiadaviek sa volia komponenty a stratégiu rozdelenia výpočtovej záťaže medzi edge a cloud.

\subsection{Výber procesorovej platformy}
Okrem základných kritérií (CPU, pamäť, periférie) by mal návrh zohľadniť:
\begin{itemize}
	\item Ekosystémový support (toolchain, debug nástroje, knižnice),
	\item Dostupnosť bezpečnostných modulov (TPM, secure element),
	\item Podporu energeticky efektívnych režimov (deep sleep, power domains),
	\item Certifikačné požiadavky a prevádzkové prostredie (teplota, vlhkosť).
\end{itemize}

\subsection{Časovanie, plánovanie a deterministické správanie}
Pri systémoch s real‑time požiadavkami je potrebné definovať časové záväzky a garantované latencie. Plánovanie úloh môže byť preemptívne (RTOS) alebo kooperatívne; výber algoritmu plánovania (fixed priority, EDF) závisí od požiadaviek na deterministickosť a jednoduchosti analýzy.

\subsection{Verifikácia a testovanie}
Okrem tradičných testovacích postupov (jednotkové testy, integračné testy, systémové testy) je dôležité zaviesť:
\begin{itemize}
	\item návrhové overenie (model checking, statická analýza),
	\item testovanie za hraníc prevádzkových podmienok (stres, teplota),
	\item overenie elektromagnetickej kompatibility (EMC/EMI),
	\item bezpečnostné testy (penetration testing, fuzzing komunikačných rozhraní).
\end{itemize}

\subsection{Regulačné a normatívne aspekty}
Pri návrhu produktov pre masové nasadenie treba zohľadniť platné normy (CE, FCC, RED) a priemyselné štandardy. Dôležité sú tiež regulácie týkajúce sa rádiového spektra a environmentálne požiadavky (RoHS).

\subsection{EMC/EMI a fyzické návrhy}
Dobre navrhnutá doska plošných spojov (PSP) a napájací reťazec minimalizujú emisné a imunitné problémy. Odporúčania zahŕňajú: oddelenie analógových a digitálnych pôd, vhodné rozmiestnenie kondenzátorov, stínovanie a sledovanie impedancií pri vysokofrekvenčných linkách.

\subsection{Napájanie a energetická efektívnosť}
Okrem komponentov napájacieho reťazca treba navrhnúť stratégiu správy energie: profilovanie spotreby, použitie adaptívnych intervalov zasielania, hibernáciu periférií a optimalizáciu algoritmov z hľadiska energie. Pre batériové systémy je kľúčové navrhnúť meranie stavu batérie a spoľahlivé mechanizmy obnovenia po vybití.

\subsection{Bezpečnostné vzory}
Odporúčané postupy zahŕňajú: bezpečné zavádzanie kľúčov pri výrobných procesoch, hardvérové zabezpečenie kľúčov (secure element), používanie overených kryptografických knižníc, zabezpečený proces OTA aktualizácií a pravidelnú bezpečnostnú údržbu.

\subsection{Nasadenie a životný cyklus produktu}
Pri plánovaní nasadenia treba definovať stratégiu aktualizácií, telemetriu pre prevádzku, monitorovanie chýb a plán údržby. Pre rozsiahle nasadenia je dôležité mať mechanizmy pre hromadný provisioning zariadení a bezpečný lifecycle management.

\subsection{Zhrnutie}
Táto kapitola poskytuje podrobný technologický a procesný rámec pre návrh, verifikáciu a prevádzku vstavaných IoT systémov. Nasledujúce kapitoly aplikujú tieto princípy v konkrétnych implementáciách na platforme ESP32 a v prostredí MicroPython.

